\section{The weight-zero scalar endpoint with one fixed ordinary factor}

This section concerns a scalar analytic identity.  It does not construct
a global Euler class or identify an arithmetic regulator.  The source
normalizations are those of Hao--Loeffler,
\emph{\(p\)-adic Rankin--Selberg \(L\)-functions in universal deformation
families and functional equations}, arXiv:2405.12611v2,
Theorem~3.5 and \S5:
\url{https://arxiv.org/pdf/2405.12611v2}.
Their printed calculation at the endpoint
\(m=k-l-2t=0\) is interpreted below through the meromorphic
Kronecker series, followed by raising.  The applied analytic comparison
is Lei--Loeffler--Zerbes, Annals~180 (2014),
Definition~4.2.1, Proposition~4.2.2, Corollary~5.2.1,
Proposition~5.2.2 and Definitions~5.3.1--5.3.3,
printed pp.~706--707 and 715--717:
\url{https://annals.math.princeton.edu/wp-content/uploads/annals-v180-n2-p06-p.pdf}.
Kato's distribution conventions are in
Ast\'erisque~295 (2004), \S\S3.6--3.10, printed pp.~138--141:
\url{https://archive.numdam.org/item/AST_2004__295__117_0.pdf}.

\begin{proposition}[AC-K32: a fixed-factor endpoint repair]
\label{AC-K32}
Let \(e\geq1\), \(p=269\), and let \(f_0\) be a fixed ordinary
cuspidal eigenform of weight \(2e+1\), new at a level prime to \(p\).
Let \(g_0\) be the weight-one cubic theta form, allowing the quadratic
compensation used above; its local representation at \(p\) is
supercuspidal of conductor exponent \(2\).
Use the tame twists, depletion, Petersson pairing and normalized
fixed-\(f_0\) projector of Hao--Loeffler's \S5.  Let \(N\), prime
to \(p\), contain the required tame levels.  Write \(\alpha\) for
the unit root of \(f_0\) and \(\beta\) for its other root.

The scalar obtained by evaluating their depleted product against
the fixed-\(f_0\) ordinary projector at
\((k,l,t)=(2e+1,1,e)\) is
\[
 2^{-2e}N^2\,\lambda_{p^2}(g_0)
 \left(\frac{p^{e+1}}{\alpha}\right)^2
 \frac{\Lambda^{[pN]}(f_0,g_0^*,e+1)}
 {(1-\beta/\alpha)(1-\beta/(p\alpha))
       \langle f_0,f_0\rangle}.
 \tag{K32.1}
\]
The notation, including the normalization of the pairing, is
exactly that of their Theorem~3.5.  This finite scalar computation
does not require residual irreducibility of \(f_0\), nor a
residually irreducible ordinary deformation family through it.
The factors \(P_p(g_0,X)\) and \(P_p(g_0^*,X)\) are \(1\);
the Fricke pseudo-eigenvalue is nonzero and both displayed
adjoint denominators are units.  In particular the scalar is
nonzero whenever the displayed imprimitive complex value is
nonzero.  No cohomological reciprocity conclusion follows from
this proposition alone.
\end{proposition}

\begin{proof}
\emph{1. Remove the weight-zero pole before specialization.}
For rational parameters \((a,b)\), use the normalized meromorphic
Kronecker series
\[
 F^{(m)}_{a,b}(\tau,s)=
 (-2\pi i)^{-m}\Gamma(s+m)\pi^{-s}y^s
 \sum_{(r,t)\ne(0,0)}
 \frac{\exp(2\pi i(rb-ta))}
 {(r\tau+t)^m|r\tau+t|^{2s}},
 \qquad y=\operatorname{Im}\tau .
 \tag{K32.2}
\]
Initially take \(\operatorname{Re}s\) sufficiently large.
The residue of \(F^{(0)}_{a,b}\) at \(s=0\) is \(-1\),
independent of \(\tau\) and of the rational parameter.
Indeed its Mellin theta integral is
\(\int_0^\infty(\Theta_{a,b}(\tau,u)-1)u^{s-1}\,du\).
Poisson summation makes the theta term exponentially small at
zero for a nontrivial twist, and \(u^{-1}\) plus an exponentially
small error for the trivial twist.  Splitting at \(u=1\)
therefore gives \(-1/s\), with an additional \(1/(s-1)\)
in the trivial case; the remaining terms are regular at zero.

In the LLZ normalization of the raising operator,
\[
 \delta_0^e F^{(0)}_{a,b}(\tau,s)
       =F^{(2e)}_{a,b}(\tau,s-e).
 \tag{K32.3}
\]
The first raising derivative kills the constant residue.
Thus the regular value at \(s=0\) exists.  The functional
equation identifies it with
\(E^{(2e)}_{a,b}(\tau,1-e)\).
Corollary~5.2.1 of LLZ shows it is nearly holomorphic;
Proposition~5.2.2 gives its positive Fourier coefficients
in the relevant one-parameter specialization:
\[
 \sum_{d\mid n}d^e(n/d)^{e-1}
       \left(\zeta^d+\zeta^{-d}\right).
 \tag{K32.4}
\]
The constant term is zero except possibly when \(e=1\);
full \(p\)-depletion removes it in that case as well.
This supplies the raised object needed in the calculation,
without postulating a holomorphic classical weight-zero
Eisenstein series.

\emph{2. Check distribution and matrix factors before taking the value.}
All finite distributions needed at this endpoint are identities
of (K32.2) on an absolute-convergence region.  For example,
for a positive integer \(D\),
\[
 \begin{split}
 \sum_{c\bmod D}F^{(m)}_{a,b+c/D}(\tau,s)
   &=D^{1-s}F^{(m)}_{a,Db}(D\tau,s),\\
 \sum_{c\bmod D}F^{(m)}_{a+c/D,b}(\tau,s)
   &=D^{1-m-s}F^{(m)}_{Da,b}(\tau/D,s),\\
 \sum_{\substack{Dx=a\\Dy=b}}F^{(m)}_{x,y}(\tau,s)
   &=D^{2-m-2s}F^{(m)}_{a,b}(\tau,s).
 \end{split}
 \tag{K32.5}
\]
The first sum filters \(r\) to multiples of \(D\); write
\(r=Du\) and use \(\operatorname{Im}(D\tau)=Dy\).
The second filters \(t=Dv\), then factors \(D\) from
\(r\tau+Dv\), using
\(\operatorname{Im}(\tau/D)=y/D\).
The third filters both coordinates.  This proves all
three powers of \(D\), including their \(s\)-dependence.

For \(\gamma=\left(\begin{smallmatrix}A&B\\C&D\end{smallmatrix}\right)
\in\mathrm{SL}_2(\mathbf Z)\), reindex
\((r',t')=(r,t)\gamma\).
The denominator and the imaginary-part factor cancel under
the weight-\(m\) slash, while the symplectic phase changes
the row parameter to \((a,b)\gamma\).
Hence
\(F_{a,b}^{(m)}|_m\gamma=F_{(a,b)\gamma}^{(m)}\)
at all \(s\).  Meromorphic continuation preserves these
identities.

For the slash convention in Hao--Loeffler,
\[
 F|_m\gamma=
 \det(\gamma)^{m-1}(c\tau+d)^{-m}F(\gamma\tau),
\]
direct differentiation gives
\[
 \delta_m^e(F|_m\gamma)
  =\det(\gamma)^{-e}
       (\delta_m^eF)|_{m+2e}\gamma,
 \qquad
 \delta_m^e B_d=d^e B_d\delta_m^e,
 \quad (B_dF)(\tau)=F(d\tau).
 \tag{K32.6}
\]
A scalar matrix \(dI\) acts by \(d^{m-2}\).
In the Fricke calculation the actual matrix identity is
\[
 B_{p^r}W_{p^{2r}}
  =(p^r I)
    \begin{pmatrix}p^r x&y\\Nz&p^r w\end{pmatrix}
    B_{p^r},
 \tag{K32.7}
\]
with the middle matrix in \(\mathrm{SL}_2(\mathbf Z)\).
Moving the raising operator past \(W_{p^{2r}}\) gives
\(p^{2re}\), and the central matrix at \(m=0\) gives
\(p^{-2r}\).  Together with the source factor
\(p^{r(m-1)}\) in the finite character sum, these are
exactly the factors in Hao--Loeffler Lemmas~5.5--5.6
at \(m=0,t=e\).
One continues (K32.5)--(K32.7), applies \(\delta_0^e\),
and then evaluates \(s=0\).  Any derivative of an
\(s\)-dependent scalar multiplies the \(\tau\)-constant
residue already killed by raising; no additional logarithmic
determinant term remains.

\emph{3. The ordinary integration by parts is an identity on actual
\(p\)-adic forms.}
Use the family of LLZ Definition~5.3.1:
\[
 E_0^{[p]}=E_{1/N}(0,-1),\qquad
 \theta^e E_0^{[p]}=E_{1/N}(e,e-1).
 \tag{K32.8}
\]
The coefficients are bounded since the depleted indices are
prime to \(p\).  LLZ Definition~5.3.3 realizes the ordinary
projection of \(E(\phi_1,\phi_2)g\) for a fixed \(g\);
it requires no ordinarity of that second factor.
On bounded \(q\)-expansions,
\(U_p\theta=p\theta U_p\), so
\[
 U_p^{n!}\theta h=p^{n!}\theta U_p^{n!}h
       \longrightarrow0.
\]
Consequently \(e_{\mathrm{ord}}\theta h=0\).
The product rule, iterated \(e\) times, gives the exact
integration-by-parts identity
\[
 \lambda_{f_0}(g\,\theta^e E_0^{[p]})
   =(-1)^e\lambda_{f_0}((\theta^e g)E_0^{[p]}).
 \tag{K32.9}
\]
This is the step used before the Petersson calculation.

The difference between full depletion and the one-variable
depletion used in Hao--Loeffler Lemma~5.4 has positive
Fourier support only at multiples of \(p\).
Its product with \(g^{[pN]}\) has no coefficient at a
multiple of \(p\), and is killed by \(U_p\), hence by
the ordinary projector.  A possible constant difference
gives a scalar multiple of \(g^{[pN]}\), killed for the
same reason.  The nearly-holomorphic comparison here is
the LLZ/Katz holomorphic-\(q\)-expansion comparison; this
does not discard nonholomorphic terms from a Petersson
integral without justification.

\emph{4. Finish the finite scalar calculation.}
The fixed-\(f_0\) projector and denominator identities
of Hao--Loeffler Lemmas~5.1 and 5.3 are unchanged.
Steps~1--3 supply the precise weight-zero replacement
in Lemmas~5.4--5.6.  Their finite Fricke, trace and
Petersson-unfolding computation through Lemma~5.8
therefore applies with \(k=2e+1,l=1,t=e\).

In Theorem~3.5, \(a=0\), \(b=2\), and the \(p\)-parts
of both nebentypus characters are trivial, including
after the quadratic compensation.  The supercuspidal
local factor has \(P_p(g_0,X)=P_p(g_0^*,X)=1\).
Substitution gives (K32.1).
Since \(v_p(\alpha)=0\), \(v_p(\beta)=2e\),
both adjoint denominator factors are units.
The Fricke operator is invertible, so its
pseudo-eigenvalue is nonzero.  The Petersson norm is
nonzero.  This proves the asserted implication from
the complex value to the scalar value.  No assertion
about a Galois cohomology class was used.
\end{proof}

\begin{remark}[Scope of the analytic repair]
The full residual-deformation hypotheses in Hao--Loeffler
are needed for their universal family theorem.
AC-K32 only replays the finite fixed-\(f_0\) calculation
with an explicit endpoint replacement, so it does not
claim that a residually reducible \(f_0\) satisfies those
universal hypotheses.  Published analytic Rankin
nonvanishing and local Euler-factor estimates may be
applied to the complex value with their precise
normalizations; they are not a proof that any arithmetic
class projects nontrivially to AC-K33.
The missing actual coefficient-class reciprocity and
the local two-dimensional regulator rank remain separate.
\end{remark}
